\subsection{Tổng quan kiến trúc ứng dụng}
\label{sec:implemented-architecture}

\subsubsection{Tổ chức theo MVVM, Repository và Model}
BroomSharing được xây dựng bằng Java và XML, sử dụng ViewModel để quản lý trạng thái màn hình và Repository để thực hiện các thao tác dữ liệu \parencite{broomsharing_source}. Cách tổ chức này vận dụng khuyến nghị phân tách giao diện với nguồn dữ liệu của Android Developers, trong đó giao diện gửi hành động xử lý và quan sát trạng thái được cung cấp bởi ViewModel \parencite{android_views,android_data}.

XML xác định bố cục; Activity tiếp nhận tương tác, gắn dữ liệu và điều hướng. Các Adapter chuyển danh sách phòng, booking hoặc hồ sơ xét duyệt thành những phần tử hiển thị. ViewModel giữ dữ liệu đang sử dụng, lựa chọn trên màn hình và kết quả thao tác qua LiveData. Khi dữ liệu thay đổi, Activity quan sát trạng thái để cập nhật nội dung, tiến trình tải hoặc thông báo lỗi. Những request tới Firestore và dịch vụ mạng được tổ chức tại Repository hoặc client tích hợp tương ứng.

Model biểu diễn dữ liệu nghiệp vụ như UserProfile, Room, Booking và ChatMessage. Các đối tượng này được sử dụng khi đọc document, truyền dữ liệu giữa các thành phần và kiểm tra điều kiện nghiệp vụ; chúng không tạo thành một dịch vụ độc lập. Hình~\ref{fig:mvvm} thể hiện quan hệ giữa các thành phần trong ứng dụng và những nguồn dữ liệu bên ngoài, sử dụng biểu tượng từ bộ Fluent UI System Icons của Microsoft \parencite{fluent_system_icons}.

\begin{figure}[htbp]
\centering
\includegraphics[width=\linewidth]{chapter5/application-architecture.pdf}
\caption{Kiến trúc tổng quan của ứng dụng BroomSharing}
\label{fig:mvvm}
\end{figure}

Firebase Authentication quản lý danh tính đăng nhập, còn Cloud Firestore lưu dữ liệu nghiệp vụ. Nhánh lưu ảnh trong hình tương ứng với Cloudinary. Nhánh bản đồ kết hợp MapLibre, MapTiler và openrouteservice. Nhánh thanh toán được triển khai qua backend Java kết nối payOS; Android không gọi API merchant bằng khóa bí mật của kênh thanh toán. Các dịch vụ này được sử dụng theo trách nhiệm riêng, trong khi dữ liệu liên kết tài khoản, phòng và lượt thuê vẫn được quản lý tại Firestore.

\subsubsection{Luồng xử lý và cập nhật trạng thái}
Với màn hình quản trị, AdminDashboardActivity gửi lựa chọn hàng đợi tới AdminViewModel. ViewModel chuyển trạng thái sang đang tải rồi gọi AdminRepository. Repository truy vấn đúng nhóm document, ánh xạ kết quả thành dữ liệu hiển thị và trả kết quả qua callback. LiveData được cập nhật sau đó để Activity hiển thị danh sách hoặc thông báo tương ứng. Khi gửi quyết định xét duyệt, giao diện khóa thao tác lặp trong lúc request chưa kết thúc và tải lại dữ liệu sau khi nhận kết quả thành công.

Đối với đặt chỗ, các điều kiện về thời gian, giá và quyền thao tác được kiểm tra lại tại tầng dữ liệu. Transaction cập nhật booking cùng lịch và bộ đếm liên quan; trạng thái thanh toán được backend Java đối soát trước khi ghi nhận. Giao diện hiển thị kết quả đã được trả về hoặc nhận từ listener, không suy ra thanh toán thành công chỉ từ việc mở QR hay quay về ứng dụng. Cơ chế transaction của Firestore cho phép các thay đổi trong cùng giao dịch được áp dụng nhất quán \parencite{firebase_transaction}.
